\documentclass[french]{book}
\usepackage{teach}
\usepackage{verbatim}
\usepackage{amsmath,amsfonts,amssymb}
\usepackage{pgfplots}
\pgfplotsset{compat=newest}
%\usepackage{geometry}
%\geometry{top=1cm, bottom=1cm, left=2cm, right=2cm}
\usepackage[utf8]{inputenc}
\usepackage[french]{babel}
\redefinecolor{thm}{title}{bgcolor}{alizarin}
\redefinecolor{prop}{title}{bgcolor}{meatbrown}
\redefinecolor{defn}{title}{bgcolor}{olivine}
\definecolor{alizarin}{rgb}{0.82, 0.1, 0.26}
\definecolor{meatbrown}{rgb}{0.9, 0.72, 0.23}
\definecolor{olivine}{rgb}{0.6, 0.73, 0.45}
\usepackage{pgf,tikz,tkz-tab}

\usepackage{mathrsfs}
\usetikzlibrary{arrows}

\begin{document}
 \setcounter{chapter}{10}
\chapter{Fonctions affines, équations de droites et systèmes}

\textit{C'est au XVIIème siècle que le mathématicien René Descartes affiche dans  l'ouvrage "La géométrie" (1637), son intension d'unir algèbre et géométrie. Il propose de réduire les problèmes de géométrie à des calculs qui se traduise bien souvent par des équations algébriques.}


\section{Définir et représenter une fonction affine}

\begin{defn}
Soit $m \in \mathbb{R}$ et $p \in \mathbb{R}$. La fonction $f$ qui à tout réel $x$ lui associe $mx+p$, s'appelle une \underline{fonction affine}. Sa représentation graphique est une droite $\Delta$.

Si $p=0$, $f$ est la \underline{fonction linéaire} de coefficient $m$, définie pour tout réel $x$ par $f(x)=mx$. $\Delta$ passe par l'origine du repère.

Si $m=0$, $f$ est une \underline{fonction constante}, définie pour tout réel $x$ par $f(x)=p$. $\Delta$ est parallèle à l'axe des abscisses.
\end{defn}

\begin{rem}
Si $f$ est définie sur $\mathbb{R}$ par $f(x)=mx+p$, les coordonnées $x$ et $y$ de tout point de la droite $\Delta$ sont liées par la relation $y=mx+p$, si bien que $\Delta= \{ (x;y)  \vert x \in \mathbb{R}, \; y=mx+p \} $ 
\end{rem}

On dit que la droite $\Delta$  a pour \underline{équation réduite} $y=mx+p$. $m$ est appelé \underline{ coefficient directeur} de $\Delta$ (ou \underline{pente} de $\Delta$ si le repère est orthornormé), $p$ est appelé \underline{ordonnée à l'origine} de $\Delta$.

\section{Proportionnalité des accroissements}


\begin{prop}
Soit $f$ une fonction affine défini sur $\mathbb{R}$. Quels que soient les réels $a$ et $b$ distincts, $\frac{f(b)-f(a)}{b-a} = m$. Le nombre $\frac{f(b)-f(a)}{b-a}$ s'appelle taux d'accroissement de $f$ entre $a$ et $b$.
\end{prop}


\section{Sens de variation d'une fonction affine}

\begin{prop}
Soit $f$ la fonction affine définie sur $\mathbb{R}$ par $f(x)=mx+p$ où $m$ et $p$ sont deux réels fixés.

\begin{itemize}
\item Si \underline{$m>0$} alors la fonction $f$ est \underline{strictement croissante} sur $\mathbb{R}$.
\item Si \underline{$m<0$} alors la fonction $f$ est \underline{strictement décroissante} sur $\mathbb{R}.$.
\item Si \underline{$m=0$} alors la fonction $f$ est \underline{constante} sur $\mathbb{R}$.
\end{itemize}
\end{prop}



\section{Vecteur directeur d'une droite}

\begin{defn}
On appelle \underline{vecteur directeur} d'une droite $\Delta$ tout vecteur $\overrightarrow{u}$ non nul tel qu'il existe deux points $A$ et $B$ de $\Delta$ vérifiant $\overrightarrow{u}=\overrightarrow{AB}$  
\end{defn}

Par conséquent, une droite possède une infinité de vecteurs directeurs mais qui ont tous la même direction, ainsi ils sont tous colinéaires entre eux.

\begin{prop}
La droite $\Delta$ passant par $A$ et de vecteur directeur $\overrightarrow{u}$ non nul est l'ensemble des points $M$ tel que $\overrightarrow{AM}$ et  $\overrightarrow{u}$ sont colinéaires.
\end{prop}

\section{Équation cartésienne et équation réduite d'une droite}

\begin{prop}
Soit $\Delta$ une droite quelconque du plan, les coordonnées $(x;y)$ de tout point de $\Delta$ sont liées par une relation de la forme $ax+by+c=0$, où $a$ \underline{et} $b$ ne peuvent pas être tous les deux nuls. 
\end{prop}

\begin{prop}
Soit $a,b$ et $c$ trois réels tels que $(a;b)\neq (0;0)$. L'ensemble $\Delta$ des points de coordonnées $(x;y)$ vérifiant $ax+by+c=0$ est une droite dont un vecteur directeur est $\overrightarrow{u} \left( \begin{array}{c}
-b \\
a \\
\end{array} \right). \\ 
$
\end{prop}

\newpage

\subsection{Établir la forme générale d'une équation de droite à l'aide du déterminant}
 
Soit un point du plan $A=(x_A,y_A)$ , un vecteur $\overrightarrow{u}\left( \begin{array}{c}
x_u\\
y_u \\
\end{array} \right) $, $\Delta$ la droite passant par $A$ dirigée par $\overrightarrow{u}$. 

On pose $M=(x,y)$ et l'on souhaite déterminer une équation en fonction de $x$ et $y$ pour traduire le fait que $M \in \Delta$. 

On a $\overrightarrow{AM} \left( \begin{array}{c}
x-x_A \\
y-y_A \\
\end{array} \right)  
$ ainsi $M \in \Delta$ \underline{si et seulement si} $\overrightarrow{AM}$ et $\overrightarrow{u}$ sont colinéaires.

Autrement dit, $M \in \Delta$ si et seulement si $det(\overrightarrow{AM},\overrightarrow{u})=0$, ce qui équivaut à l'équation suivante:
 $$y_u(x-x_A) -x_u(y-y_A)=0$$
ou encore, 
$$y_u x - x_u y + x_u y_A - x_A y_u = 0 .$$
 
\begin{defn}
L'équation $ax+by+c=0$ est appelée \underline{une équation cartésienne de la droite $\Delta$}
\vspace{3mm}
\end{defn}

Certaines droites peuvent représenter des fonctions affines. Si une droite $\Delta$ a pour équation $(E): ax+by+c=0$, avec $b \neq 0$, alors $\Delta$ n'est pas parallèle à l'axe des ordonnées. L'équation $(E)$ équivaut à $y=- \frac{a}{b} x - \frac{c}{b}$ on retrouve alors \underline{l'équation réduite}  $y=mx+p$ en posant $m=-\frac{a}{b}$ et $p=-\frac{c}{b}$.\\

La droite $\Delta$ représente graphiquement alors la fonction $f(x)=mx+p$, $m$ est appelé le \underline{ le coefficient directeur} ou \underline{pente} de la droite $\Delta$.

\begin{prop}
Soit $\Delta$ une droite d'équation réduite $y=mx+p$. Un vecteur directeur de $\Delta$ est $\overrightarrow{u} \left( \begin{array}{c}
1 \\
m \\
\end{array} \right). \\ 
$
\end{prop}

\subsection*{Le cas des droites parallèles à l'axe des ordonnées}
\begin{prop}
Une droite $\Delta$ du plan est parrallèle à l'axe des ordonnées si et seulement s'il existe $k \in \mathbb{R}$ tel que $x-k=0$ (ou encore $x=k$) est une équation cartésienne de $\Delta$. 
\end{prop}

\section{Droites parallèles}
\begin{prop}
Soient deux droites $\Delta$ et $\Delta^{'}$ dont une équation cartésienne respective des deux droites sont $(E)\; ax+by+c=0$ et $(E')\; a'x+b'y+c'=0$. \\

Alors $\Delta$ et $\Delta^{'}$ sont parallèles si l'une des conditions suivantes est vérifiée.

\begin{itemize}
\item Un vecteur directeur de $\Delta$ est colinéaire à un vecteur directeur de $\Delta^{'}$.
\item Leurs coefficients directeurs sont égaux
\item $ab'-ba'=0$
\end{itemize}

\end{prop}

\section{Système et intersection}

On se donne $a,b,c,a',b',c' \in \mathbb{R}$ des nombres réels.

\begin{defn}
On appelle \underline{système} de deux équations linéaire à deux inconnues la donnée simultanée de deux équations du premier degré à deux inconnues: 

$$
\left \{
\begin{array}{rcl}
ax+by&=&c \\
a'x+b'y&=&c'
\end{array}
\right. 
$$

Un couple $(x;y)$ est dit \underline{solution} du système s'il vérifie \underline{à la fois les deux égalités du système}.

\end{defn}

Soient deux droites $\Delta$ et $\Delta^{'}$ dont une équation cartésienne respective des deux droites sont $(E)\; ax+by-c=0$ et $(E')\; a'x+b'y-c'=0$. Ainsi on a $(a;b) \neq (0;0)$ et $(a',b') \neq (0;0)$.

Alors le système: 

$$
\left \{
\begin{array}{rcl}
ax+by&=&c \\
a'x+b'y&=&c'
\end{array}
\right. 
$$ 

Permet d'étudier l'intersection des deux droites  $\Delta$ et $\Delta^{'}$. 
Effectivement on sait que l'équation de la ligne 1 (notée $(1)$) équivaut à $(E)\; ax+by-c=0$. Donc si un couple $(x;y)$ vérifie l'équation $(1)$ alors c'est un point de $\Delta$. De même l'équation de la ligne 2 (notée $(2)$) équivaut à  $(E')\; a'x+b'y-c'=0$. Donc si un couple $(x;y)$ vérifie l'équation $(2)$ alors c'est un point de $\Delta^{'}$. \\

Donc si un couple $(x;y)$ est solution du système alors c'est un point de $\Delta$ \underline{et} de $\Delta^{'}$.

On pourra cependant distinguer deux cas.

\subsection*{Le cas où $ ab'-b'a \neq 0$}
$\Delta$ et $\Delta^{'}$ sont sécantes et le système a une \underline{unique} solution.
\subsection*{Le cas où $ ab'-b'a = 0$}
\subsubsection*{Soit $\Delta$ et $\Delta^{'}$ sont strictement parallèles}
Alors le système n'a \underline{aucune} solution.
\subsubsection*{Soit $\Delta$ et $\Delta^{'}$ sont confondus}
Alors le système a une \underline{infinité} de solution.

\begin{prop}
Le système admet une \underline{unique solution}  si et seulement si $ab'-a'b \neq 0$.
\end{prop}

\end{document}
